\documentclass[aspectratio=169, xcolor=dvipsnames]{beamer}

% --- Theme Setup ---
\usetheme{Madrid}
\usecolortheme{default}

% --- Packages ---
\usepackage{graphicx}
\usepackage{xcolor}
\usepackage{booktabs}
\usepackage{hyperref}
\usepackage{adjustbox}
\usepackage{amssymb}
\usepackage{amsmath}
\usepackage{listings}
\usepackage{caption}

\lstset{
  basicstyle=\ttfamily\small,
  keywordstyle=\color{blue}\bfseries,
  stringstyle=\color{red},
  showstringspaces=false,
  breaklines=true
}

% --- Title Information ---
\title{Module 11}
\subtitle{SQL Examples}
\author{Milav Dabgar}
\institute{www.milav.in}
\date{\today}

\begin{document}

% Title Slide
\begin{frame}
    \titlepage
\end{frame}

\begin{frame}{Week Recap}
    \begin{itemize}[<+->]
        \item Basic notions of Relational Database Models
        \item Attributes and their types
        \item Mathematical structure of relational model
        \item Schema and Instance
        \item Keys, primary as well as foreign
        \item Relational algebra with operators
        \item Relational query language
        \item DDL (Data Definition)
        \item DML (Basic Query Structure)
        \item Detailed understanding of basic query structure
        \item Set operations, null values, and aggregation
    \end{itemize}
\end{frame}

\begin{frame}{Module Objectives}
    \begin{block}{Goals}
        \begin{itemize}[<+->]
            \item To recap various basic SQL features through example workout
        \end{itemize}
    \end{block}
\end{frame}

\begin{frame}{Module Outline}
    \begin{itemize}[<+->]
        \item Examples of basic SQL
    \end{itemize}
\end{frame}

\section{SELECT}
\begin{frame}
  \vfill
  \centering
  \begin{beamercolorbox}[sep=8pt,center,shadow=true,rounded=true]{title}
    \usebeamerfont{title}SELECT\par%
  \end{beamercolorbox}
  \vfill
\end{frame}

\begin{frame}[fragile]{Select distinct}
    \begin{exampleblock}{Example}
        \begin{itemize}[<+->]
            \item From the \textbf{classroom} relation in the figure, find the names of buildings in which every individual classroom has capacity less than 100 (removing the duplicates).
        \end{itemize}
    \end{exampleblock}
    \begin{columns}
        \begin{column}{0.45\textwidth}
            \textbf{classroom} relation
            \begin{center}
                \begin{adjustbox}{max width=\textwidth}
                \begin{tabular}{|l|l|l|}
                \hline
                \textbf{building} & \textbf{room\_number} & \textbf{capacity} \\ \hline
                Packard & 101 & 500 \\ \hline
                Painter & 514 & 10 \\ \hline
                Taylor & 3128 & 70 \\ \hline
                Watson & 100 & 30 \\ \hline
                Watson & 120 & 50 \\ \hline
                \end{tabular}
                \end{adjustbox}
            \end{center}
        \end{column}
        \begin{column}{0.55\textwidth}
            \textbf{Query:}
            \begin{lstlisting}[language=SQL]
select distinct building 
from classroom 
where capacity < 100;
            \end{lstlisting}
            \textbf{Output:}
            \begin{center}
                \begin{adjustbox}{max width=\textwidth}
                \begin{tabular}{|l|}
                \hline
                \textbf{building} \\ \hline
                Painter \\ \hline
                Taylor \\ \hline
                Watson \\ \hline
                \end{tabular}
                \end{adjustbox}
            \end{center}
        \end{column}
    \end{columns}
\end{frame}

\begin{frame}[fragile]{Select all}
    \begin{exampleblock}{Example}
        \begin{itemize}[<+->]
            \item From the \textbf{classroom} relation in the figure, find the names of buildings in which every individual classroom has capacity less than 100 (without removing the duplicates).
        \end{itemize}
    \end{exampleblock}
    \begin{columns}
        \begin{column}{0.45\textwidth}
            \textbf{classroom} relation
            \begin{center}
                \begin{adjustbox}{max width=\textwidth}
                \begin{tabular}{|l|l|l|}
                \hline
                \textbf{building} & \textbf{room\_number} & \textbf{capacity} \\ \hline
                Packard & 101 & 500 \\ \hline
                Painter & 514 & 10 \\ \hline
                Taylor & 3128 & 70 \\ \hline
                Watson & 100 & 30 \\ \hline
                Watson & 120 & 50 \\ \hline
                \end{tabular}
                \end{adjustbox}
            \end{center}
        \end{column}
        \begin{column}{0.55\textwidth}
            \textbf{Query:}
            \begin{lstlisting}[language=SQL]
select all building 
from classroom 
where capacity < 100;
            \end{lstlisting}
            \textbf{Output:}
            \begin{center}
                \begin{adjustbox}{max width=\textwidth}
                \begin{tabular}{|l|}
                \hline
                \textbf{building} \\ \hline
                Painter \\ \hline
                Taylor \\ \hline
                Watson \\ \hline
                Watson \\ \hline
                \end{tabular}
                \end{adjustbox}
            \end{center}
            \small Note that duplicate retention is the default and hence it is a common practice to skip \texttt{all} immediately after \texttt{select}.
        \end{column}
    \end{columns}
\end{frame}

\section{Cartesian Product / AS}
\begin{frame}
  \vfill
  \centering
  \begin{beamercolorbox}[sep=8pt,center,shadow=true,rounded=true]{title}
    \usebeamerfont{title}Cartesian Product / AS\par%
  \end{beamercolorbox}
  \vfill
\end{frame}

\begin{frame}[fragile]{Cartesian Product}
    \begin{block}{Overview}
        \begin{itemize}[<+->]
            \item Find the list of all students of departments which have a budget < \$0.1 million
            \begin{lstlisting}[language=SQL]
select name, budget 
from student, department 
where student.dept_name = department.dept_name 
  and budget < 100000;
            \end{lstlisting}
            \item The above query first generates every possible student-department pair, which is the Cartesian product of \texttt{student} and \texttt{department}. Then, it filters all the rows with \texttt{student.dept\_name = department.dept\_name} and \texttt{budget < 100000}.
            \item The common attribute \texttt{dept\_name} in the resulting table are renamed using the relation name (\texttt{student.dept\_name} and \texttt{department.dept\_name})
        \end{itemize}
    \end{block}
\end{frame}

\begin{frame}[fragile]{Rename AS Operation}
    \begin{block}{Query}
        \begin{itemize}[<+->]
            \item The same query in the previous slide can be framed by renaming the tables as shown below.
            \begin{lstlisting}[language=SQL]
select S.name as studentname, budget as deptbudget 
from student as S, department as D 
where S.dept_name = D.dept_name 
  and budget < 100000;
            \end{lstlisting}
            \item The above query renames the relation \texttt{student} as \texttt{S} and the relation \texttt{department} as \texttt{D}
            \item It also displays the attribute \texttt{name} as \texttt{StudentName} and \texttt{budget} as \texttt{DeptBudget}.
            \item Note that the \texttt{budget} attribute does not have any prefix because it occurs only in the \texttt{department} relation.
        \end{itemize}
    \end{block}

    \textbf{Output:}
    \vspace{-0.5em}
    \begin{center}
        \begin{adjustbox}{max width=0.5\textwidth}
        \begin{tabular}{|l|l|}
        \hline
        \textbf{studentname} & \textbf{deptbudget} \\ \hline
        Brandt & 50000 \\ \hline
        Peltier & 70000 \\ \hline
        Levy & 70000 \\ \hline
        Sanchez & 80000 \\ \hline
        Snow & 70000 \\ \hline
        Aoi & 85000 \\ \hline
        Bourikas & 85000 \\ \hline
        Tanaka & 90000 \\ \hline
        \end{tabular}
        \end{adjustbox}
    \end{center}
\end{frame}

\section{WHERE: AND / OR}
\begin{frame}
  \vfill
  \centering
  \begin{beamercolorbox}[sep=8pt,center,shadow=true,rounded=true]{title}
    \usebeamerfont{title}WHERE: AND / OR\par%
  \end{beamercolorbox}
  \vfill
\end{frame}

\begin{frame}[fragile]{Where: AND and OR}
    \begin{exampleblock}{Example}
        \begin{itemize}[<+->]
            \item From the \texttt{instructor} and \texttt{department} relations in the figure, find out the names of all instructors whose department is 'Finance' or whose department is in any of the following buildings: 'Watson', 'Taylor'.
            \item \textbf{Query:}
            \begin{lstlisting}[language=SQL, basicstyle=\ttfamily\scriptsize]
select name 
from instructor I, department D 
where D.dept_name = I.dept_name 
  and (I.dept_name = 'Finance' or building in ('Watson','Taylor'));
            \end{lstlisting}
            \item \textbf{Output:}
        \end{itemize}
    \end{exampleblock}
    \begin{center}
        \begin{adjustbox}{max width=0.5\textwidth}
        \begin{tabular}{|l|}
        \hline
        \textbf{name} \\ \hline
        Srinivasan \\ \hline
        Wu \\ \hline
        Einstein \\ \hline
        Gold \\ \hline
        Katz \\ \hline
        Singh \\ \hline
        Crick \\ \hline
        Brandt \\ \hline
        Kim \\ \hline
        \end{tabular}
        \end{adjustbox}
    \end{center}
\end{frame}

\begin{frame}{Instructor and Department Relations}
    \textbf{department}
    \vspace{-0.5em}
    \begin{center}
        \begin{adjustbox}{max width=\textwidth}
        \begin{tabular}{|l|l|l|}
        \hline
        \textbf{dept\_name} & \textbf{building} & \textbf{budget} \\ \hline
        Biology & Watson & 90000 \\ \hline
        Comp. Sci. & Taylor & 100000 \\ \hline
        Elec. Eng & Taylor & 85000 \\ \hline
        Finance & Painter & 120000 \\ \hline
        History & Painter & 50000 \\ \hline
        Music & Packard & 80000 \\ \hline
        Physics & Watson & 70000 \\ \hline
        \end{tabular}
        \end{adjustbox}
    \end{center}
    
    \textbf{instructor}
    \vspace{-0.5em}
    \begin{center}
        \begin{adjustbox}{max width=\textwidth}
        \begin{tabular}{|l|l|l|l|}
        \hline
        \textbf{ID} & \textbf{name} & \textbf{dept\_name} & \textbf{salary} \\ \hline
        10101 & Srinivasan & Comp. Sci. & 65000 \\ \hline
        12121 & Wu & Finance & 90000 \\ \hline
        15151 & Mozart & Music & 40000 \\ \hline
        22222 & Einstein & Physics & 95000 \\ \hline
        32343 & El Said & History & 60000 \\ \hline
        33456 & Gold & Physics & 87000 \\ \hline
        45565 & Katz & Comp. Sci. & 75000 \\ \hline
        58583 & Califieri & History & 62000 \\ \hline
        76543 & Singh & Finance & 80000 \\ \hline
        76766 & Crick & Biology & 72000 \\ \hline
        83821 & Brandt & Comp. Sci. & 92000 \\ \hline
        98345 & Kim & Elec. Eng & 80000 \\ \hline
        \end{tabular}
        \end{adjustbox}
    \end{center}
\end{frame}

\section{String}
\begin{frame}
  \vfill
  \centering
  \begin{beamercolorbox}[sep=8pt,center,shadow=true,rounded=true]{title}
    \usebeamerfont{title}String\par%
  \end{beamercolorbox}
  \vfill
\end{frame}

\begin{frame}[fragile]{String Operations}
    \begin{block}{Pattern Matching}
        \begin{itemize}[<+->]
            \item From the \texttt{course} relation in the figure, find the titles of all courses whose course id has three alphabets indicating the department.
            \item \textbf{Query:}
            \begin{lstlisting}[language=SQL]
select title 
from course 
where course_id like '___-%';
            \end{lstlisting}
            \item \textbf{Output:}
        \end{itemize}
    \end{block}
    \begin{center}
        \begin{adjustbox}{max width=0.5\textwidth}
        \begin{tabular}{|l|}
        \hline
        \textbf{title} \\ \hline
        Intro. to Biology \\ \hline
        Genetics \\ \hline
        Computational Biology \\ \hline
        Investment Banking \\ \hline
        World History \\ \hline
        Physical Principles \\ \hline
        \end{tabular}
        \end{adjustbox}
    \end{center}
    
    \uncover<4->{
        \small The course id of each department has either 2 or 3 alphabets in the beginning, followed by a hyphen and then followed by a 3-digit number. The above query returns the names of those departments that have 3 alphabets in the beginning.
    }
\end{frame}

\begin{frame}{course relation}
    \begin{center}
        \begin{adjustbox}{max width=\textwidth}
        \begin{tabular}{|l|l|l|l|}
        \hline
        \textbf{course\_id} & \textbf{title} & \textbf{dept\_name} & \textbf{credits} \\ \hline
        BIO-101 & Intro. to Biology & Biology & 4 \\ \hline
        BIO-301 & Genetics & Biology & 4 \\ \hline
        BIO-399 & Computational Biology & Biology & 3 \\ \hline
        CS-101 & Intro. to Computer Science & Comp. Sci. & 4 \\ \hline
        CS-190 & Game Design & Comp. Sci. & 4 \\ \hline
        CS-315 & Robotics & Comp. Sci. & 3 \\ \hline
        CS-319 & Image Processing & Comp. Sci. & 3 \\ \hline
        CS-347 & Database System Concepts & Comp. Sci. & 3 \\ \hline
        EE-181 & Intro. to Digital Systems & Elec. Eng. & 3 \\ \hline
        FIN-201 & Investment Banking & Finance & 3 \\ \hline
        HIS-351 & World History & History & 3 \\ \hline
        MU-199 & Music Video Production & Music & 3 \\ \hline
        PHY-101 & Physical Principles & Physics & 4 \\ \hline
        \end{tabular}
        \end{adjustbox}
    \end{center}
\end{frame}

\section{ORDER BY}
\begin{frame}
  \vfill
  \centering
  \begin{beamercolorbox}[sep=8pt,center,shadow=true,rounded=true]{title}
    \usebeamerfont{title}ORDER BY\par%
  \end{beamercolorbox}
  \vfill
\end{frame}

\begin{frame}[fragile]{Order By}
    \begin{block}{Overview}
        \begin{itemize}[<+->]
            \item From the \texttt{student} relation in the figure, obtain the list of all students in alphabetic order of departments and within each department, in decreasing order of total credits.
            \item The list is first sorted in alphabetic order of \texttt{dept\_name}.
            \item Within each dept, it is sorted in decreasing order of \texttt{tot\_cred}.
            \item \textbf{Query:}
            \begin{lstlisting}[language=SQL]
select name, dept_name, tot_cred 
from student 
order by dept_name ASC, tot_cred DESC;
            \end{lstlisting}
        \end{itemize}
    \end{block}
\end{frame}

\begin{frame}[fragile]{Order By - student relation and output}
    \begin{columns}
        \begin{column}{0.5\textwidth}
            \textbf{student}
            \begin{center}
                \begin{adjustbox}{max width=\textwidth, max height=0.7\textheight}
                \begin{tabular}{|l|l|l|l|}
                \hline
                \textbf{ID} & \textbf{name} & \textbf{dept\_name} & \textbf{tot\_cred} \\ \hline
                00128 & Zhang & Comp. Sci. & 102 \\ \hline
                12345 & Shankar & Comp. Sci. & 32 \\ \hline
                19991 & Brandt & History & 80 \\ \hline
                23121 & Chavez & Finance & 110 \\ \hline
                44553 & Peltier & Physics & 56 \\ \hline
                45678 & Levy & Physics & 46 \\ \hline
                54321 & Williams & Comp. Sci. & 54 \\ \hline
                55739 & Sanchez & Music & 38 \\ \hline
                70557 & Snow & Physics & 0 \\ \hline
                76543 & Brown & Comp. Sci. & 58 \\ \hline
                76653 & Aoi & Elec. Eng. & 60 \\ \hline
                98765 & Bourikas & Elec. Eng. & 98 \\ \hline
                98988 & Tanaka & Biology & 120 \\ \hline
                \end{tabular}
                \end{adjustbox}
            \end{center}
        \end{column}
        \begin{column}{0.5\textwidth}
            \textbf{Output}
            \begin{center}
                \begin{adjustbox}{max width=\textwidth, max height=0.7\textheight}
                \begin{tabular}{|l|l|l|}
                \hline
                \textbf{name} & \textbf{dept\_name} & \textbf{tot\_cred} \\ \hline
                Tanaka & Biology & 120 \\ \hline
                Zhang & Comp. Sci. & 102 \\ \hline
                Brown & Comp. Sci. & 58 \\ \hline
                Williams & Comp. Sci. & 54 \\ \hline
                Shankar & Comp. Sci. & 32 \\ \hline
                Bourikas & Elec. Eng. & 98 \\ \hline
                Aoi & Elec. Eng. & 60 \\ \hline
                Chavez & Finance & 110 \\ \hline
                Brandt & History & 80 \\ \hline
                Sanchez & Music & 38 \\ \hline
                Peltier & Physics & 56 \\ \hline
                Levy & Physics & 46 \\ \hline
                Snow & Physics & 0 \\ \hline
                \end{tabular}
                \end{adjustbox}
            \end{center}
        \end{column}
    \end{columns}
\end{frame}

\section{IN}
\begin{frame}
  \vfill
  \centering
  \begin{beamercolorbox}[sep=8pt,center,shadow=true,rounded=true]{title}
    \usebeamerfont{title}IN\par%
  \end{beamercolorbox}
  \vfill
\end{frame}

\begin{frame}[fragile]{In Operator}
    \begin{block}{Overview}
        \begin{itemize}[<+->]
            \item From the \texttt{teaches} relation in the figure, find the IDs of all courses taught in the Fall or Spring of 2018.
            \item \emph{Note: We can use \texttt{distinct} to remove duplicates.}
            \item \textbf{Query:}
            \begin{lstlisting}[language=SQL]
select course_id 
from teaches 
where semester in ('Fall','Spring') and year = 2018;
            \end{lstlisting}
        \end{itemize}
    \end{block}
\end{frame}

\begin{frame}{In Operator - teaches relation and output}
    \begin{columns}
        \begin{column}{0.6\textwidth}
            \textbf{teaches}
            \begin{center}
                \begin{adjustbox}{max width=\textwidth, max height=0.7\textheight}
                \begin{tabular}{|l|l|l|l|}
                \hline
                \textbf{ID} & \textbf{course\_id} & \textbf{semester} & \textbf{year} \\ \hline
                10101 & CS-101 & Fall & 2017 \\ \hline
                10101 & CS-315 & Spring & 2018 \\ \hline
                10101 & CS-347 & Fall & 2017 \\ \hline
                12121 & FIN-201 & Spring & 2018 \\ \hline
                15151 & MU-199 & Spring & 2018 \\ \hline
                22222 & PHY-101 & Fall & 2017 \\ \hline
                32343 & HIS-351 & Spring & 2018 \\ \hline
                45565 & CS-101 & Spring & 2018 \\ \hline
                45565 & CS-319 & Spring & 2018 \\ \hline
                76766 & BIO-101 & Summer & 2017 \\ \hline
                76766 & BIO-301 & Summer & 2018 \\ \hline
                83821 & CS-190 & Spring & 2017 \\ \hline
                83821 & CS-190 & Spring & 2017 \\ \hline
                83821 & CS-319 & Spring & 2018 \\ \hline
                98345 & EE-181 & Spring & 2017 \\ \hline
                \end{tabular}
                \end{adjustbox}
            \end{center}
        \end{column}
        \begin{column}{0.4\textwidth}
            \textbf{Output}
            \begin{center}
                \begin{adjustbox}{max width=\textwidth}
                \begin{tabular}{|l|}
                \hline
                \textbf{course\_id} \\ \hline
                CS-315 \\ \hline
                FIN-201 \\ \hline
                MU-199 \\ \hline
                HIS-351 \\ \hline
                CS-101 \\ \hline
                CS-319 \\ \hline
                CS-319 \\ \hline
                \end{tabular}
                \end{adjustbox}
            \end{center}
        \end{column}
    \end{columns}
\end{frame}

\section{Set}
\begin{frame}
  \vfill
  \centering
  \begin{beamercolorbox}[sep=8pt,center,shadow=true,rounded=true]{title}
    \usebeamerfont{title}Set\par%
  \end{beamercolorbox}
  \vfill
\end{frame}

\begin{frame}[fragile]{Set Operations: UNION}
    \begin{block}{Overview}
        \begin{itemize}[<+->]
            \item For the same question in the previous slide, we can find the solution using \texttt{union} operator as follows.
            \item \textbf{Query:}
            \begin{lstlisting}[language=SQL]
select course_id from teaches 
where semester='Fall' and year=2018 

union 

select course_id from teaches 
where semester='Spring' and year=2018;
            \end{lstlisting}
            \item Note that \texttt{union} removes all duplicates. If we use \texttt{union all} instead of \texttt{union}, we get the same set of tuples as in previous slide.
        \end{itemize}
    \end{block}
\end{frame}

\begin{frame}[fragile]{Set Operations (2): INTERSECT}
    \begin{exampleblock}{Example}
        \begin{itemize}[<+->]
            \item From the \texttt{instructor} relation, find the names of all instructors who taught in either the Computer Science department or the Finance department and whose salary is < 80000.
            \item \textbf{Query:}
            \begin{lstlisting}[language=SQL]
select name 
from instructor 
where dept_name in ('Comp. Sci.','Finance') 

intersect 

select name 
from instructor 
where salary < 80000;
            \end{lstlisting}
            \item Note that the same can be achieved using the query:
            \begin{lstlisting}[language=SQL, basicstyle=\ttfamily\scriptsize]
select name from instructor 
where dept_name in ('Comp. Sci.', 'Finance') and salary < 80000;
            \end{lstlisting}
        \end{itemize}
    \end{exampleblock}
\end{frame}

\begin{frame}[fragile]{Set Operations (3): EXCEPT}
    \begin{exampleblock}{Example}
        \begin{itemize}[<+->]
            \item From the \texttt{instructor} relation, find the names of all instructors who taught in either the Computer Science department or the Finance department and whose salary is either $\ge 90000$ or $\le 70000$.
            \item \textbf{Query:}
            \begin{lstlisting}[language=SQL, basicstyle=\ttfamily\scriptsize]
select name 
from instructor 
where dept_name in ('Comp. Sci.','Finance') 

except 

select name 
from instructor 
where salary < 90000 and salary > 70000;
            \end{lstlisting}
            \item Note that the same can be achieved using the query given below:
            \begin{lstlisting}[language=SQL, basicstyle=\ttfamily\scriptsize]
select name from instructor 
where dept_name in ('Comp. Sci.', 'Finance') 
  and (salary >= 90000 or salary <= 70000);
            \end{lstlisting}
        \end{itemize}
    \end{exampleblock}
\end{frame}

\section{Aggregation}
\begin{frame}
  \vfill
  \centering
  \begin{beamercolorbox}[sep=8pt,center,shadow=true,rounded=true]{title}
    \usebeamerfont{title}Aggregation\par%
  \end{beamercolorbox}
  \vfill
\end{frame}

\begin{frame}[fragile]{Aggregate functions: AVG}
    \begin{columns}
        \begin{column}{0.6\textwidth}
            \begin{block}{Average}
                \begin{itemize}
                    \item From the \texttt{classroom} relation, find the names and the average capacity of each building whose average capacity is greater than 25.
                    \item \textbf{Query:}
                    \begin{lstlisting}[language=SQL]
select building, avg(capacity) 
from classroom 
group by building 
having avg(capacity) > 25;
                    \end{lstlisting}
                \end{itemize}
            \end{block}
        \end{column}
        \begin{column}{0.4\textwidth}
            \textbf{Output:}
            \begin{center}
                \begin{adjustbox}{max width=\textwidth}
                \begin{tabular}{|l|l|}
                \hline
                \textbf{building} & \textbf{avg} \\ \hline
                Taylor & 70 \\ \hline
                Packard & 500 \\ \hline
                Watson & 40 \\ \hline
                \end{tabular}
                \end{adjustbox}
            \end{center}
        \end{column}
    \end{columns}
\end{frame}

\begin{frame}[fragile]{Aggregate functions (2): MIN}
    \begin{block}{Minimum}
        \begin{itemize}[<+->]
            \item From the \texttt{instructor} relation, find the least salary drawn by any instructor among all the instructors.
            \item \textbf{Query:}
            \begin{lstlisting}[language=SQL]
select min(salary) as least_salary 
from instructor;
            \end{lstlisting}
            \item \textbf{Output:}
        \end{itemize}
    \end{block}
    \begin{center}
        \begin{adjustbox}{max width=0.3\textwidth}
        \begin{tabular}{|l|}
        \hline
        \textbf{least\_salary} \\ \hline
        40000 \\ \hline
        \end{tabular}
        \end{adjustbox}
    \end{center}
\end{frame}

\begin{frame}[fragile]{Aggregate functions (3): MAX}
    \begin{block}{Maximum}
        \begin{itemize}[<+->]
            \item From the \texttt{student} relation, find the maximum credits obtained by any student among all the students.
            \item \textbf{Query:}
            \begin{lstlisting}[language=SQL]
select max(tot_cred) as max_credits 
from student;
            \end{lstlisting}
            \item \textbf{Output:}
        \end{itemize}
    \end{block}
    \begin{center}
        \begin{adjustbox}{max width=0.3\textwidth}
        \begin{tabular}{|l|}
        \hline
        \textbf{max\_credits} \\ \hline
        120 \\ \hline
        \end{tabular}
        \end{adjustbox}
    \end{center}
\end{frame}

\begin{frame}[fragile]{Aggregate functions (4): COUNT}
    \begin{block}{Count}
        \begin{itemize}[<+->]
            \item From the \texttt{section} relation, find the number of courses run in each building.
            \item \textbf{Query:}
            \begin{lstlisting}[language=SQL]
select building, count(course_id) as course_count 
from section 
group by building;
            \end{lstlisting}
            \item \textbf{Output:}
        \end{itemize}
    \end{block}
    \begin{center}
        \begin{adjustbox}{max width=0.4\textwidth}
        \begin{tabular}{|l|l|}
        \hline
        \textbf{building} & \textbf{course\_count} \\ \hline
        Taylor & 5 \\ \hline
        Packard & 4 \\ \hline
        Painter & 3 \\ \hline
        Watson & 3 \\ \hline
        \end{tabular}
        \end{adjustbox}
    \end{center}
\end{frame}

\begin{frame}[fragile]{Aggregate functions (5): SUM}
    \begin{block}{Sum}
        \begin{itemize}[<+->]
            \item From the \texttt{course} relation, find the total credits offered by each department.
            \item \textbf{Query:}
            \begin{lstlisting}[language=SQL]
select dept_name, sum(credits) as sum_credits 
from course 
group by dept_name;
            \end{lstlisting}
            \item \textbf{Output:}
        \end{itemize}
    \end{block}
    \begin{center}
        \begin{adjustbox}{max width=0.4\textwidth}
        \begin{tabular}{|l|l|}
        \hline
        \textbf{dept\_name} & \textbf{sum\_credits} \\ \hline
        Finance & 3 \\ \hline
        History & 3 \\ \hline
        Physics & 4 \\ \hline
        Music & 3 \\ \hline
        Comp. Sci. & 17 \\ \hline
        Biology & 11 \\ \hline
        Elec. Eng. & 3 \\ \hline
        \end{tabular}
        \end{adjustbox}
    \end{center}
\end{frame}

\begin{frame}{Module Summary}
    \begin{block}{Summary}
        \begin{itemize}[<+->]
            \item SQL Examples have been practiced for
            \begin{itemize}
                \item[$\triangleright$] Select
                \item[$\triangleright$] Cartesian Product / \texttt{as}
                \item[$\triangleright$] Where: \texttt{and} / \texttt{or}
                \item[$\triangleright$] String Matching
                \item[$\triangleright$] \texttt{Order by}
                \item[$\triangleright$] \texttt{in}
                \item[$\triangleright$] Set Operations: \texttt{union}, \texttt{intersect}, \texttt{except}
                \item[$\triangleright$] Aggregate Functions: \texttt{avg}, \texttt{min}, \texttt{max}, \texttt{count}, \texttt{sum}
            \end{itemize}
        \end{itemize}
    \end{block}
    \vspace{2em}
    \begin{center}
    \small \textit{Slides used in this presentation are borrowed from \url{http://db-book.com/} with kind permission of the authors.} \\
    \small \textit{Edited and new slides are marked with 'PPD'.}
    \end{center}
\end{frame}

\end{document}
